\definecolor{ocre}{RGB}{243,102,25} % Define the orange color used for highlighting throughout the book
\definecolor{aliceblue}{rgb}{0.94, 0.97, 1.0}	% Defintion colback
\definecolor{byzantium}{rgb}{0.44, 0.16, 0.39}  % Defintion colframe

\definecolor{richelectricblue}{rgb}{0.03, 0.57, 0.82} 	% Example colback
\definecolor{powderblue}{rgb}{0.69, 0.88, 0.9} 	% Example colframe

\definecolor{cornsilk}{rgb}{1.0, 0.97, 0.86}
\definecolor{scarlet}{rgb}{1.0, 0.13, 0.0}
%----------------------------------------------------------------------------------------
%	THEOREM STYLES
%----------------------------------------------------------------------------------------

\usepackage{amsmath,amsfonts,amssymb,amsthm} % For math equations, theorems, symbols, etc

\newcommand{\intoo}[2]{\mathopen{]}#1\,;#2\mathclose{[}}
\newcommand{\ud}{\mathop{\mathrm{{}d}}\mathopen{}}
\newcommand{\intff}[2]{\mathopen{[}#1\,;#2\mathclose{]}}

\renewcommand{\qedsymbol}{$\blacksquare$}% Optional qed square
\usepackage{tcolorbox}
\tcbuselibrary{theorems}

% May need to add breakable, but perhaps each theorem should be continuous, and not split up.
\newtcbtheorem[number within=section]{theorem}{Theorem}%
{colback=orange!5,colframe=black,fonttitle=\bfseries}{th}

\newtcbtheorem[number within=section]{definition}{Definition}%
{colback=aliceblue,colframe=byzantium,fonttitle=\bfseries}{th}

\newtcbtheorem[number within=section]{example}{Example}%
{colback=powderblue,colframe=richelectricblue,fonttitle=\bfseries}{th}

\newtcbtheorem[number within=section]{corollary}{Corollary}%
{colback=cornsilk,colframe=scarlet!50!black,fonttitle=\bfseries}{th}
%----------------------------------------------------------------------------------------
%	REMARK ENVIRONMENT
%----------------------------------------------------------------------------------------

\newenvironment{remark}{\par\vspace{10pt}\small % Vertical white space above the remark and smaller font size
	\begin{list}{}{
			\leftmargin=35pt % Indentation on the left
			\rightmargin=25pt}\item\ignorespaces % Indentation on the right
		\makebox[-2.5pt]{\begin{tikzpicture}[overlay]
			\node[draw=ocre!60,line width=1pt,circle,fill=ocre!25,font=\sffamily\bfseries,inner sep=2pt,outer sep=0pt] at (-15pt,0pt){\textcolor{ocre}{R}};\end{tikzpicture}} % Orange R in a circle
		\advance\baselineskip -1pt}{\end{list}\vskip5pt} % Tighter line spacing and white space after remark
